\noindent
{\Large \textbf{\section*{Заключение}}} \par
\addcontentsline{toc}{section}{Заключение} \par

В ходе выполнения лабораторной работы разработана программа на языке C, определяющая характеристики кэш-памяти компьютера по косвенным измерениям времени обращения к памяти. Для процессора Apple M4 получены следующие значения: размер строки кэша - 128 байт, размер кэша первого уровня - 128 КБ, размер кэша второго уровня - 16 МБ, ассоциативность кэша первого уровня - 8.

\vspace{8pt}
Реализовано:
\vspace{8pt}
\begin{enumerate}
    \item Измерение среднего времени обращения к памяти по цепочке зависимых обращений с прогревом и выбором медианы по нескольким замерам.
    \item Определение размера строки по зависимости времени от шага обхода, с автоматическим выбором шага, на котором время возрастает сильнее всего.
    \item Определение размеров уровней кэша по зависимости времени от размера массива, с хаотичным порядком обхода строк для исключения влияния префетчера.
    \item Определение ассоциативности по числу ячеек, конфликтующих за один набор кэша.
\end{enumerate}

% \vspace{8pt}
% \textbf{Достоинства реализованного решения.} Программа компактна, не использует системные вызовы получения информации о кэше, а найденный размер строки передается в следующий эксперимент, так что он не задается заранее. Полученные значения совпадают с параметрами процессора, выдаваемыми операционной системой.

% \vspace{8pt}
% \textbf{Недостатки и возможности для улучшения.}
% \begin{enumerate}
%     \item Размеры массивов перебираются степенями двойки, поэтому границы кэшей определяются с точностью до удвоения (например, размер L1 лежит в промежутке между 128 и 256 КБ, а за размер принимается последний быстрый). Точность можно повысить более мелким шагом по размеру.
%     \item Третий уровень кэша не выявлен: на использованном процессоре после кэша второго уровня отдельного кэша третьего уровня, размер которого можно определить таким способом, обнаружено не было.
%     \item Ассоциативность надежно определяется только для первого уровня, так как индекс набора в кэше второго уровня вычисляется по физическому адресу, а программа управляет лишь виртуальными адресами.
%     \item В macOS нет возможности привязать поток к ядру, поэтому результат может зависеть от того, на каком ядре выполняется программа.
% \end{enumerate}

% \vspace{8pt}
% \textbf{Масштабируемость решения.} Объем памяти и диапазон шагов задаются константами в программе, поэтому метод применим на других компьютерах без изменения алгоритма, но для процессоров с большими кэшами необходимо увеличить максимальный размер массива. Время работы программы растет линейно с числом обращений и числом замеров.
